\section{Experimental Design}

\subsection{Scenarios}

The evaluation focuses on correctness, decryption success, failure behavior, timing, and storage overhead. Fifteen scenarios were selected to cover both the threshold property and the Reed-Solomon correction bound. Each scenario was executed for 25 independent trials, giving 375 total runs.

Table~\ref{tab:scenarios} lists the evaluated configurations. The ``available'' count $m$ denotes the number of shares presented to the decoder, while $e$ denotes the number of shares intentionally corrupted before recovery. Scenarios marked as failure are expected to fail either because $m < t$ or because the redundancy condition $m \geq t + 2e$ is not satisfied.

\begin{table*}[!t]
\caption{Experimental Scenarios and Observed Outcomes Across 25 Trials Per Scenario}
\label{tab:scenarios}
\centering
\scriptsize
\begin{tabular}{p{2.7in}cccccc}
\toprule
\textbf{Scenario} & \textbf{$t$} & \textbf{$n$} & \textbf{$m$} & \textbf{$e$} & \textbf{Observed Success Rate} & \textbf{Mean Recovery (ms)} \\
\midrule
Normal recovery & 3 & 5 & 5 & 0 & 100\% & 0.1419 \\
Minimum valid recovery & 3 & 5 & 3 & 0 & 100\% & 0.0736 \\
Insufficient shares & 3 & 5 & 2 & 0 & 0\% & 0.0017 \\
One corrupted share, enough redundancy & 3 & 5 & 5 & 1 & 100\% & 0.1781 \\
One corrupted share, not enough redundancy & 3 & 5 & 4 & 1 & 0\% & 0.0402 \\
Larger normal recovery & 4 & 7 & 7 & 0 & 100\% & 0.2937 \\
Larger minimum recovery & 4 & 7 & 4 & 0 & 100\% & 0.1375 \\
Larger insufficient shares & 4 & 7 & 3 & 0 & 0\% & 0.0018 \\
Correct one corrupted share & 4 & 7 & 6 & 1 & 100\% & 0.2397 \\
Correct two corrupted shares & 4 & 8 & 8 & 2 & 100\% & 0.4009 \\
Large threshold normal & 5 & 10 & 10 & 0 & 100\% & 0.6451 \\
Large threshold minimum & 5 & 10 & 5 & 0 & 100\% & 0.1254 \\
Large threshold insufficient & 5 & 10 & 4 & 0 & 0\% & 0.0018 \\
Large threshold, one corrupted & 5 & 10 & 7 & 1 & 100\% & 0.2841 \\
Large threshold, two corrupted & 5 & 10 & 9 & 2 & 100\% & 0.4787 \\
\bottomrule
\end{tabular}
\end{table*}

\subsection{Measurement Procedure}

Each trial follows the same procedure. First, a fresh AES-256 key is generated and used to encrypt a benchmark plaintext with AES-GCM. Second, the key is split into shares under the specified $(t,n)$ configuration. Third, a random subset of $m$ shares is selected. Fourth, the scenario optionally corrupts $e$ of those shares by modifying their $y$ values. Finally, the recovery stage attempts to reconstruct the key and decrypt the ciphertext. Success is recorded only when the recovered key exactly matches the original key and AES-GCM verification succeeds.

Timing is measured separately for share generation and for recovery. Storage overhead is computed as the total serialized share size divided by the original 32-byte AES key size. Because the serialization format is fixed-width, the overhead depends only on the number of generated shares.

\subsection{Evaluation Metrics}

Five metrics are used in the evaluation. First, \textit{recovery success} records whether the robust decoder returns a candidate secret without violating the algebraic constraints. Second, \textit{key correctness} checks whether the recovered 32-byte value matches the original AES key exactly. Third, \textit{decryption success} checks whether AES-GCM accepts the recovered key and reproduces the original plaintext. Fourth, \textit{recovery latency} measures the wall-clock time of the reconstruction stage. Fifth, \textit{storage overhead} measures the total share footprint relative to the 32-byte secret.

These metrics distinguish between algebraic recovery and application-level recovery. In principle a decoder might output a field element even when it is not the original key. By checking both exact key equality and authenticated decryption, the evaluation rejects such cases. This distinction is important because a key-recovery system is useful only when the recovered secret is operationally valid for the protected ciphertext.
